\begin{center}
\LARGE
Computer Algebra with Data of Limited Accuracy 

\end{center}

\begin{center}
\Large
Hans J. Stetter
\end{center}

\noindent
Date:  July 19th (Friday)\\
Time:  14:00-14:30\\

\begin{center}
\large
Abstract
\end{center}

In Scientific Computing, some of the numerical data in problems (e.g. coefficients in polynomials)
have limited accuracy. The application of standard computer algebra systems to such problems is
unsatisfactory for various reasons. An adaptation of polynomial computer algebra for data with
limited accuracy requires - first of all - a change of a fundamental paradigm: Data spaces have to
be furnished with a topology. It turns out that a number of classical algebraic concepts (like g.c.d.,
Groebner bases, etc.) are n o t continuous w.r.t. their data in situations which are smooth w.r.t
these data. 

We will display various such representation singularities and exhibit their potentially disastrous
consequences for the respective algorithms. We will also point out how this difficulty may be
overcome by a suitable modification of the conceptual approaches. The resulting algorithms are
well within the scope of computer algebra. 

